%% Chapter 3, Section 3.4: Discussion and Bibliography
%% A First Course in Monte Carlo Methods — Sanz-Alonso & Al-Ghattas

\section{Discussion and Bibliography}
\label{sec:3.4}

The book chapter \cite{105} is a classic reference on the Monte Carlo method and variance
reduction techniques. Modern textbooks include \cite{14, 179, 90, 130, 145, 193, 95, 15,
173, 194}. Some of these books emphasize specific applications, such as finance
\cite{90, 234, 116}, statistics \cite{14, 193, 179}, scientific computing \cite{145},
statistical physics \cite{173}, stochastic simulation \cite{95}, or computer vision,
machine learning, and artificial intelligence \cite{15}. For a textbook on quantum Monte
Carlo methods with applications in chemistry, we refer to \cite{99}. The textbook
\cite{194} introduces Monte Carlo methods through examples coded in R, while \cite{130}
includes numerous worked examples using MATLAB\@. Finally, \cite{40} provides an accessible
overview of Monte Carlo methods aimed at applied mathematicians.

In Section~\ref{sec:3.1} we invoked the central limit theorem to derive an asymptotic
confidence interval for classical Monte Carlo integration. Non-asymptotic analyses are also
possible, and we refer to \cite{95}, Chapter~3, for further details. Empirical process theory
provides another perspective on Monte Carlo methods, see e.g.\ \cite{231}, Chapter~8, for a
gentle introduction.

The lecture notes \cite{9} compare Monte Carlo and importance sampling through examples.
Importance sampling was developed as a variance reduction technique in the early 1950's
\cite{121, 120}. Theorem~\ref{thm:3.3}, which shows how to optimally choose the proposal
density for a given test function $h$ and target density $f$, was already shown in
\cite{121}. For recent methodological developments, see \cite{140, 178, 146, 222, 122}. The
book \cite{49} contains a modern presentation of importance sampling in a general framework.
A review of importance sampling, from the perspective of filtering and sequential importance
resampling, can be found in \cite{3}; the proofs in this chapter closely follow the
presentation in that paper. For a review of adaptive importance sampling techniques, see
\cite{39}.

The key role of the second moment of the weight function $w$ has long been realized
\cite{144, 186}, and it is known to be asymptotically linked to the effective sample size
\cite{128, 129, 144}. The question of how to choose a proposal distribution that leads to
small value of said second moment has been widely studied, and we refer to \cite{142} and
references therein. Similar to \cite{46}, Theorem~\ref{thm:3.5} quantifies the estimation
error in terms of the $\chi^2$-divergence between target and proposal; recent complementary
analysis of importance sampling in \cite{45} utilizes the Kullback--Leibler divergence.
Necessary sample size results for importance sampling in terms of several divergences between
target and proposal were established in \cite{207, 209}. A review of useful distances
between probability measures can be found in \cite{86}. Effective sample size based on
discrepancy measures are studied in \cite{156}. The paper \cite{71} overviews existing
notions of effective sample size for importance sampling algorithms and proposes new ones.

Variance reduction techniques are covered in most standard textbooks on Monte Carlo methods,
see e.g.\ \cite{14}, Chapter~5, and \cite{193}, Chapter~4. The monograph \cite{95} considers
more advanced variance reduction techniques. In particular, we refer to \cite{95}, Part~III,
for variance reduction techniques for simulation of stochastic differential equations.
